\chapter{Feb.~26 --- Differentials, Part 2}

\section{Tangent Sheaf}

\begin{prop}\label{prop:differentials-affine-variety}
  We have the following:
  \begin{enumerate}
    \item For $X \subseteq \Affine^n$
      closed, we have
      \[\Omega_X(X) = \frac{\bigoplus_{i = 1}^n A dx_i}{(df_1, \dots, df_r)},\]
      where $A = \OO_X(X) = k[x_1, \dots, x_n] / I(X)$
      and $I(X) = (f_1, \dots, f_r) \le k[x_1, \dots, x_n]$. 
    \item For $x \in X$,
      we have $(\Omega_X)_x = \Omega_{\OO_{X, x} / k}$.
  \end{enumerate}
\end{prop}

\begin{proof}
  (1) This was the discussion in
  Example \ref{example:affine-variety-differentials}.

  (2) We may assume that $X$ is affine.
  Let $I := I_X(\{x\}) \le \OO_X(X)$.
  Then
  \[
    (\Omega_X)_x
    \cong (\widetilde{\Omega_{\OO_X(X) / k}})_x
    \cong (\Omega_{\OO_X(X) / k})_I
    = \Omega_{\OO_{X, x} / k},
  \]
  which is the desired claim.
\end{proof}

\begin{prop}
  For a variety $X$ and $x \in X$,
  \begin{enumerate}
    \item
      $T_{X, x}^\vee \cong \Omega_X(x)$,
      in particular
      $\dim \Omega_X(x) = \dim_k T_{X, x}$;
    \item $X$ is smooth if and only if
      $\Omega_X$ is locally free.
  \end{enumerate}
  In particular, when
  (2) holds, $T_X = \Omega_X^\vee$ is a
  locally free sheaf.
\end{prop}

\begin{proof}
  (1) We will only show the
  dimension statement.
  As both are local statements,
  we may assume that $x = 0 \in X \subseteq \Affine^n$
  where $X$
  is closed. Write $A = \OO_X(X)$
  and $I_X = (f_1, \dots, f_r) \le k[x_1, \dots, x_n]$.
  Then
  \[
    \Omega_X(0)
    = \frac{\Omega_X(X)}{(\overline{x}_1, \dots, \overline{x}_n) \Omega_X(X)}
    = \frac{\bigoplus_{i = 1}^n k dx_i}{(\sum_j (\partial f_i / \partial x_j)(0) dx_j)}.
  \]
  by Proposition
  \ref{prop:differentials-affine-variety}.
  As $T_{X, 0} = V(f_1^{\mathrm{lin}}, \dots, f_r^{\mathrm{lin}})$,
  we see that $\dim \Omega_X(0) = \dim T_{X, 0}$.

  (2) We may assume that $X$ is connected.
  If $X$ is smooth, then $X$ is irreducible,
  so
  \[
    n := \dim X
    = \codim_X\{x\} = \dim T_{X, x}
    = \dim \Omega_X(x)
  \]
  for all $x \in X$. Thus a homework
  problem implies that $\Omega_X$ is
  free.

  Conversely, assume $\Omega_X$
  is free. Set $n = \rank \Omega_X$,
  then for $x \in X^{\mathrm{sm}}$,
  we have $\codim_X\{x\} = \dim \Omega_X(x) = n$.
  Now $X^{\mathrm{sm}}$ is dense in $X$,
  so $X$ is of pure dimension $n$.
  So for all $x \in X$, we have
  \[
    n = \dim X = \codim_X\{x\}
    \le \dim T_{x, x} = \dim \Omega_X(x)
    = n,
  \]
  hence $\dim T_{X, x} = n$, so $X$
  is smooth.
\end{proof}

\begin{definition}
  For a variety $X$, its \emph{tangent
  sheaf} is $T_X := \Omega_X^\vee$.
\end{definition}

\begin{remark}
  If $X$ is smooth, then $T_X$ is
  locally free.
\end{remark}

\begin{conjecture}[Zariski-Lipman]
  If $T_X$ is locally free, then
  $X$ is smooth.
\end{conjecture}

\begin{remark}
  In general, $\mathcal{F}$ being
  locally free does not imply that
  $\mathcal{F}^\vee$ is locally free.
\end{remark}
